\chapter*{Remerciements}

Je tiens à remercier en premier lieu mon encadrante pédagogique, \textbf{Mme Loubna Karbil}, pour son suivi tout au long de ce stage, ses conseils, sa disponibilité et ses orientations méthodologiques. Son accompagnement m'a permis de mieux structurer mon travail, d'approfondir ma réflexion et de mener à bien les différentes étapes de ce projet.

Mes sincères remerciements s'adressent également à mon encadrante au sein de \textbf{S. WATRANS}, \textbf{Mme Sara Zaqi}, pour m'avoir accueilli au sein de l'entreprise et plus particulièrement au sein du département informatique. Je la remercie pour la confiance qu'elle m'a accordée en me confiant le projet \textbf{SmartTrans-AI}, ainsi que pour sa disponibilité, ses conseils et les échanges qui m'ont permis de mieux comprendre les besoins et les problématiques liés au domaine du transport et de la gestion de flotte.

Je souhaite également exprimer ma gratitude à l'ensemble des collaborateurs de \textbf{S. WATRANS} pour leur accueil, leur disponibilité et leur bienveillance durant toute la période de mon stage. Leur collaboration et leurs explications m'ont permis de découvrir davantage le fonctionnement d'une entreprise et de mieux comprendre l'importance des outils informatiques dans la gestion quotidienne de ses activités.

Ce stage a constitué pour moi une expérience particulièrement enrichissante, aussi bien sur le plan technique que professionnel. Il m'a permis de mettre en pratique les connaissances acquises durant ma formation, notamment dans les domaines du développement logiciel, des bases de données, de la conception des systèmes et de l'intelligence artificielle. Il m'a également permis de développer mon autonomie, mon esprit d'analyse, ma capacité à résoudre des problèmes et mon aptitude à travailler sur un projet répondant à des besoins réels.

Je remercie également l'ensemble du corps professoral de l'\textbf{ENSIASD de Taroudant} pour la qualité de l'enseignement dispensé et pour les connaissances et compétences acquises au cours de ma formation. Ces acquis ont constitué une base importante pour la réalisation de ce stage et du projet qui m'a été confié.

Je tiens aussi à remercier mes camarades de promotion pour leur soutien, leurs échanges d'idées et les discussions techniques qui ont contribué à rendre cette expérience plus enrichissante.

Enfin, j'adresse mes plus sincères remerciements à ma famille pour son soutien constant, ses encouragements, sa patience et sa confiance. Leur accompagnement m'a donné la motivation nécessaire pour avancer et donner le meilleur de moi-même tout au long de cette période.

À toutes les personnes qui ont contribué, directement ou indirectement, à la réussite de ce stage et à la réalisation de ce projet, je leur adresse mes sincères remerciements et ma profonde gratitude.
